\section{Arutelu} \label{arutelu}

Käesolev peatükk arutleb eksperimendi tulemuste üle laiemas kontekstis, käsitleb TI kasutamise praktilisi aspekte, piiranguid ja annab soovitused.

\subsection{TI tugevused lõputöö kirjutamisel} \label{subsec:tugevused}

Eksperimendi tulemused kinnitavad, et kaasaegsed TI mudelid koos agentsete tööriistadega suudavad pakkuda märkimisväärset abi lõputöö kirjutamise protsessis. Selle katse jooksul demonstreeris Claude Opus 4.6 + Cowork kombinatsioon eriti veenvalt, et tipptasemel mudel koos agentse tööriistaga suudab genereerida tervikliku 61-leheküljelise eestikeelse lõputöö ka ühe kasutajaviibega ilma iteratiivse suunamiseta. Peamised tugevused jagunevad järgmiselt.

\textbf{Struktureerimine ja kavandamine.} TI mudelid on eriti tugevad teksti struktureerimisel. Kõik neli testitud kombinatsiooni (arvestades GPT-5.4 Codex-laienduse kaudu saadud tulemust) suutsid luua loogilise ja akadeemilistele nõuetele vastava peatükkide jaotuse ainuüksi ühe kasutajaviibega. Eriti kasulik on TI kasutamine töö algfaasis, kui üliõpilane vajab abi sisukorra ja peatükkide ülesehituse kavandamisel. Mudelid suudavad vaid minimaalse juhise põhjal luua detailse sisukorra koos soovitusliku leheküljemahuga iga peatüki jaoks.

\textbf{Mustandi loomine.} TI suudab genereerida tekstimustandeid, mis on piisavalt hea kvaliteediga, et neid edasi töödelda. See kiirendab kirjutamisprotsessi oluliselt, kuna üliõpilane ei pea alustama tühjalt lehelt, vaid saab toimetada ja täiendada olemasolevat teksti. Esialgse kasutuskogemuse põhjal võib TI mustandi kasutamine säästa hinnanguliselt 50--80\% kirjutamisele kuluvast ajast, kuigi täpse ajavõidu mõõtmine vajab edaspidi süstemaatilisi kasutajakatseid.

\textbf{Keeleline toimetamine.} Kõik neli mudelit suudavad efektiivselt teostada keelelist toimetamist -- parandada grammatikavigu, parandada lauseehitust ja ühtlustada stiili. See on eriti väärtuslik eesti keele kontekstis, kus spetsiifilisi akadeemilise toimetamise teenuseid on vähe.

\textbf{LaTeX-vormindamine.} Oluline praktiline eelis on mudelite võimekus genereerida teksti otse LaTeX-vormingus, mis hõlmab korrektset viitamissüntaksit, tabeleid, jooniste viiteid ja matemaatilisi valemeid. See säästab üliõpilasele märkimisväärset vormistusaega. Võrdluseks, Microsoft Wordi puhul peab üliõpilane vormistuse -- stiilid, ristviited, pealkirjade numeratsioon ja bibliograafia -- enamasti käsitsi seadistama, kuna TI mudelid ei saa \texttt{.docx} faili sisemist struktuuri samamoodi manipuleerida nagu tekstipõhist \LaTeX{} lähtekoodi. See teeb Wordi TI-toetatud akadeemilise kirjutamise kontekstis oluliselt aeganõudvamaks ja veaohtlikumaks valikuks.

\subsection{TI puudujäägid ja piirangud} \label{subsec:puudujaagid}

Vaatamata edusammudele on TI mudelitel lõputöö kirjutamise kontekstis endiselt olulised puudujäägid. Tulemuste tõlgendamisel tuleb eriti arvestada kahte metoodilist piirangut: eksperimendis võrreldi \emph{mudeli ja agentse tööriista kombinatsioone}, mitte mudeleid isoleeritult, ning iga kombinatsiooni kohta tehti vaid üks jooks. Seetõttu näitavad tulemused konkreetsete töövoogude käitumist selles katses, mitte lõplikku mudelite üldist paremusjärjestust. Süstemaatiline hinnang nõuaks iga kombinatsiooni kohta mitut kordusjooksu ning suuremat valimit.

\textbf{Viidete usaldusväärsus.} Kõige tõsisem probleem on viidete hallutsineerimine. Kõigi mudelite genereeritud lõputöödes esines olematuid allikaid. Gemini 3.1 Pro hallutsineeris viiteid ca 25\% juhtudest, GPT-5.4 ca 20\%, Claude Sonnet 4.6 ca 20\% ja Claude Opus 4.6 ca 15\% juhtudest. Kuigi Opus 4.6 näitas madalaimat hallutsinatsioonimäära, on ka 15\% liiga kõrge akadeemilise teksti jaoks. Need tulemused on kooskõlas varasemate uurimustega. Chelli et al. \cite{chelli2024hallucination} tuvastasid GPT-4 hallutsinatsioonimääraks 28,6\%, mis on lähedane meie GPT-5.4 tulemusele. See viitab sellele, et vaatamata mudelite arengule pole viidete hallutsineerimise probleem täielikult lahendatud. Mugaanyi et al. \cite{mugaanyi2024citation} leidsid samuti, et DOI-de täpsus on valdkonniti väga erinev (loodusteadustes 32,7\% vs humanitaarteadustes 8,5\%), mis viitab sellele, et arvutiteaduse valdkonnas võivad tulemused olla mõnevõrra paremad. Kim et al. \cite{kim2025geographic} on näidanud, et hallutsinatsioonid on eriti levinud väiksemate keelte ja madalama sissetulekuga riikide kontekstis, mis on asjakohane ka eesti keele puhul. See tähendab, et iga TI-genereeritud viide tuleb inimese poolt kontrollida. Viidete automaatne genereerimine on hetkel ebapiisava kvaliteediga ja üliõpilane peab ise allikaid otsima ning viiteid kontrollima.

\textbf{Sisuline sügavus.} TI-genereeritud tekstid kipuvad olema pealiskaudsed. Mudelid kordavad üldtuntud seisukohti, kuid ei suuda luua originaalset analüüsi ega teha uudseid järeldusi. See on eriti problemaatiline tulemuste arutelu puhul, kus eeldatakse, et autor demonstreerib kriitilist mõtlemist ja seob tulemused laiema kontekstiga kokku.

\textbf{Eesti keele tehniline terminoloogia.} Kuigi mudelite eesti keele oskus on oluliselt paranenud, esineb endiselt probleeme tehnilise terminoloogiaga. Mudelid kasutavad kohati ebaühtlast terminoloogiat (vahetavad eesti- ja ingliskeelsete terminite vahel) ning mõnikord loovad valesid termineid. Näiteks genereeriti juurdunud termini ``süvaõpe'' asemel ``süvaõppimine'' või kasutati sõna ``treenima'' asemel ``koolitama'', mis ei vasta arvutiteaduse valdkonnas juurdunud eesti\-keelsele terminoloogiale. Barbu et al. \cite{barbu2025estonian} on samuti täheldanud, et keelemudelid vajavad eesti keele spetsiifika jaoks täiendavat kohandamist, mis viitab laiemale probleemile väikeste keelte tehnilistel kasutusjuhtudel.

\textbf{Kontekstiakna piirangud.} Kuigi kontekstiaknad on oluliselt suurenenud, on need endiselt piiratud. 25-leheküljelise bakalaureusetöö tekst (ca 80\,000--100\,000 tähemärki ehk ca 40\,000--50\,000 tokenit) mahub küll kontekstiaknasse, kuid kui lisada juhendid, näidistööd ja varasemalt genereeritud osad, on kontekstiakna efektiivne kasutamine endiselt kriitiline.

\textbf{Juhiste järgimise ebaühtlus ning liidese ja agentsete tööriistade mõju.} Peatükis~\ref{eksperimendid} kirjeldatud tulemused näitavad, et sama mudel võib eri liidestes käituda sisuliselt erinevalt. GPT-5.4 vestlusliidese keeldumine ja Codex-laienduse kaudu õnnestunud failipõhine genereerimine viitavad sellele, et väljundit ei määra ainult alusmudel, vaid ka süsteemijuhised, ohutuspoliitika ja töövoo ülesehitus. Bianchi et al. \cite{bianchi2023safety} kirjeldatud ``liialdatud ohutuskäitumine'' aitab seda nähtust tõlgendada: liides võib uurimusliku ülesande tõlgendada akadeemilise väärkasutuse riskina. Agentne tööriist ei tähenda seejuures ohutuspiirangutest möödumist, vaid teistsugust ülesande konteksti, kus töö jaotub väiksemateks faili- ja kompileerimispõhisteks sammudeks.

Sama oluline on eristada vormilist veenvust sisulisest usaldusväärsusest. Gemini 3.1 Pro + Antigravity töö näitas kõige selgemalt, et korrektselt vormistatud ja statistiliste tabelitega tekst võib sisaldada täielikult fabritseeritud empiirikat. Opus 4.6 + Cowork andis selles katses tugevama ja tehniliselt sügavama väljundi, kuid ka seal tuli kontrollimata testitulemusi ja viiteid käsitleda riskina. Seega ei piisa agentse tööriista edukast \LaTeX{}-vormistusest: sisulise vastutuse, andmete olemasolu ja viidete kontroll peab jääma inimesele.

\textbf{Hindamismetoodika piirangud.} Oluline piirang käesolevas uurimuses on asjaolu, et hindamise viis läbi ainult üks hindaja (töö autor). See tähendab, et hindajatevahelise kooskõla mõõtmine ei olnud võimalik ja subjektiivsus on paratamatult suurem kui mitme sõltumatu hindaja kaasamisel. Subjektiivsuse vähendamiseks kasutati iga teksti korduvat hindamist eri aegadel, kuid see ei asenda sõltumatuid hindajaid. Eriti subjektiivsetes kriteeriumides (akadeemiline stiil, praktiline väärtus) võivad ühe hindaja hinnangud kalduda süstemaatiliselt ühes suunas. Edaspidistes uurimustes tuleks tulemuste kinnitamiseks kaasata mitu sõltumatut hindajat, sealhulgas valdkonna õppejõud ja teised üliõpilased.

\subsection{Hinnakujundus ja majanduslik teostatavus} \label{subsec:hinnakujundus}

Tabelis~\ref{tab:hinnavordlus} on esitatud hinnanguline kuluarvestus 25-leheküljelise bakalaureusetöö genereerimiseks iga mudeliga.

\begin{table}[ht]
\centering
\caption{Bakalaureusetöö genereerimise hinnanguline kuluarvestus}
\label{tab:hinnavordlus}
\begin{tblr}{
    width=\linewidth,
    colspec={X[1.6,l] X[1,c] X[1,c] X[1,c] X[1,c]},
    rows={font=\small},
    colsep=3pt,
    row{1}={font=\bfseries\small},
    hlines,
    vlines,
    row{even}={bg=colorCellHighlight},
}
Kulukomponent & Gemini 3.1 Pro & GPT-5.4 & Claude Sonnet 4.6 & Claude Opus 4.6 \\
Sisendtokenid (ca 2M) & \$4,00--8,00 & \$5,00 & \$6,00 & \$30,00 \\
Väljundtokenid (ca 300K) & \$3,60--5,40 & \$4,50 & \$4,50 & \$22,50 \\
Agentse tööriista sisend/väljundkorrad & mitmekordne & mitmekordne & mitmekordne & mitmekordne \\
\textbf{Eeldatav hind} & \textbf{\$23--38} & \textbf{\$25--32} & \textbf{\$32--42} & \textbf{\$160--220} \\
\end{tblr}
\end{table}

Hinnad on esitatud USA dollarites vastavalt mudelite pakkujate 2026. aasta märtsi ametlikele hinnakirjadele \cite{openai2026pricing, anthropic2026pricing, google2026gemini_pricing}; euros väljendatuna on summad sõltuvad kasutatavast valuutakursist.

Kuutellimused muudavad kulupildi mõnevõrra teistsuguseks: ChatGPT Plus ja Claude Pro maksid vaadeldud perioodil \$20/kuu ning Google AI Studio võimaldas piiratud tasuta kasutust. Praktikas võib üliõpilane teha põhilised katsetused 1--2 kuu tellimuse hinnaga, kui mahupiiranguid ei ületata.

Kokkuvõttes on TI kasutamise majanduslik kulu Sonnet-klassi kombinatsioonidega mõistlik, samas kui Opus-klassi mudeli kasutamine on märksa kallim ja vajab teadlikku kvaliteedi-kulu kaalumist. Ka kallimad TI-variandid jäävad siiski üldjuhul odavamaks kui professionaalse toimetaja teenus, mis võib bakalaureusetöö puhul maksta mitusada eurot.

\subsection{Eesti keele korrektuur ja toimetamine} \label{subsec:korrektuur}

Eraldi väärib tähelepanu TI mudelite kasutamine eestikeelse teksti korrektuuri ja toimetamise tööriistana. See on üks valdkondi, kus TI pakub kõige suuremat praktilist väärtust.

Järgmised hinnangud põhinevad autori mitteformaalsetel katsetel, mis jäid peaeksperimendi raamistikust välja ning mille täpsem metoodika vajab edaspidi süstemaatilist vormistamist. Need on esitatud suunavate tähelepanekutena.

\textbf{Grammatikakontroll.} Kõik neli mudelit suudavad tuvastada ja parandada enamiku grammatikavigadest eesti keeles. Claude Opus 4.6 oli kõige täpsem, tuvastades enamiku sisestatud vigadest; Sonnet 4.6 ja GPT-5.4 tuvastasid enamiku, kuid jätsid osa keerulisemaid vigu märkamata; Gemini 3.1 Pro tulemus oli nõrgim, jättes tuvastamata ka lihtsamaid morfoloogiavigu.

\textbf{Stiilikontroll.} Mudelid suudavad hinnata teksti vastavust akadeemilisele stiilile ja pakkuda parandusettepanekuid. See hõlmab liiga vaba keele tuvastamist, passiivi kasutamise soovitusi ja liigsete korduste märkamist. Claude Sonnet 4.6 oli selles osas kõige põhjalikum, pakkudes detailsemaid selgitusi paranduste kohta.

\textbf{Terminoloogiakontroll.} Mudelid suudavad kontrollida tehnilise terminoloogia ühtlust ja korrektsust, kuid siin on täpsus oluliselt madalam kui grammatikakontrollis. Terminoloogilist kontrolli ei saa täielikult TI-le delegeerida ja valdkonnaspetsiifiline inimekspertiis on endiselt vajalik.

\textbf{Soovitus.} TI-põhine korrektuur on kõige efektiivsem, kui seda kasutatakse kahes etapis:
\begin{enumerate}
    \item Esmalt laseb üliõpilane TI-l teksti üle vaadata ja parandusettepanekud teha;
    \item Seejärel vaatab üliõpilane parandused üle ja otsustab, millised neist aktsepteerida;
    \item Lõpuks saab üliõpilane teha käsitsi täiendavaid parandusi, mida TI ei tuvastanud ja vajadusel alustada esimesest punktist uuesti.
\end{enumerate}

Sellist kolmeetapilist protsessi kasutades saab oluliselt parandada teksti keelelist kvaliteeti, ilma et tekiks oht TI soovituste pimesi järgimiseks.

\subsection{Soovitused TI tööriista valikuks} \label{subsec:soovitused}

Eksperimendi tulemuste ja praktiliste kogemuste põhjal saab anda järgmised soovitused.

\textbf{Selle katse kõrgeim kvaliteet: Claude Opus 4.6 koos Cowork'iga.} Käesolev eksperiment näitas, et Claude Opus 4.6 + Cowork kombinatsioon saavutas selles katses kõrgeima kvaliteedi eestikeelse akadeemilise teksti genereerimisel. Kombinatsioon ületas teisi keelelise korrektsuse, akadeemilise stiili ja sisulise sügavuse poolest. Peamine puudus on oluliselt kõrgem hind. Kuna iga kombinatsiooni kohta tehti vaid üks jooks ja keelemudelid on mitteterministlikud, tuleb seda soovitust võtta eksploratiivse viitena, mitte lõpliku kvaliteedipingerea kinnitusena.

\textbf{Parim hinna-kvaliteedi suhe: Claude Sonnet 4.6 koos Cowork'iga.} Claude Sonnet 4.6 + Cowork kombinatsioon on kõige sobivam valik piiratud eelarvega üliõpilase jaoks. Kombinatsioon paistis silma akadeemilise stiili, juhiste järgimise ja teksti struktureerimise poolest ning on oluliselt soodsam kui Opus. Pikendatud mõtlemisrežiim parandab keerulisemate peatükkide kvaliteeti.

\textbf{Lühikeste lõikude jaoks: GPT-5.4.} GPT-5.4 sobib hästi lühikeste tekstilõikude genereerimiseks ja keelelise korrektuuri jaoks. Kui soov on genereerida tervikdokumenti, siis ChatGPT vestlusliidese süsteemijuhised ja ohutuspoliitika võivad selle blokeerida või suunata ingliskeelsele väljundile, samas kui ChatGPT Codex laiendus Visual Studio Code'is andis käesolevas katses sama viipega tervikliku eestikeelse dokumendi (vt peatükk~\ref{subsec:codex}).

\textbf{Ideaalne töövoog.} Käesoleva töö kogemus näitab, et kõige parema tulemuse andis Cowork'i agentne tööriist koos Opus või Sonnet mudeliga.

\subsection{Tööriistad, mida vältida} \label{subsec:valtida}

Lisaks peaeksperimendis osalenud nelja mudeli võrdlusele testis autor lühemate mitteformaalsete katsete raames ka mõnda teist tööriista, mida lõputöö kirjutamise abivahendina ei saa soovitada, vaatamata nende laialdasele kättesaadavusele üliõpilastele. Järgmised hinnangud ei põhine käesoleva uurimuse formaalsel võrdlusel, vaid autori praktilisel kasutuskogemusel ning nende arendajate avalikul dokumentatsioonil.

\textbf{Perplexity ei sobi mahukate ülesannete jaoks.} Perplexity on TI-põhine otsingumootor, mis pakub juurdepääsu erinevatele mudelitele (sh Claude, GPT ja teised). Kuigi Perplexity sobib hästi üksikute küsimuste esitamiseks ja lühikeste vastuste saamiseks, ilmnes testimisel oluline puudus: Perplexity vahendab mudeleid oma teenusekonfiguratsiooni ja päringuloogika kaudu, mistõttu kasutatav kontekst ei kattu alati mudelite natiivsete piiridega. Perplexity ametlik API-dokumentatsioon eristab otsingukonteksti mudeli kontekstiaknast ning näitab, et vahendatud mudelite piirangud sõltuvad konkreetsest Perplexity teenusest ja mudelikonfiguratsioonist \cite{perplexity2026pricing, perplexity2026models}. Autori testimisel ilmnes veebiliideses praktiliselt väiksem otsese sisendi piirang; DataStudiosi ülevaade kirjeldab selleks vaikimisi ca 8\,000 tokenit päringu kohta \cite{datastudios2025perplexity}, mis on oluliselt vähem kui samade mudelite natiivsed kontekstiaknad (nt Claude'i või Gemini 1M tokenit). See tähendab, et eraldi lühikeste promptidena saab küll normaalselt töötada. Mahukamate tekstide genereerimisel või pikema vestluse konteksti hoidmisel (nt tervikliku peatüki kirjutamine koos eelnevate osade kontekstiga) jääb kontekstiaknast aga puudu. Lisaks on täheldatud, et pikema konteksti puhul hakkab Perplexity varasemaid vestluse osi kärpima või kompresseerima, mille tagajärjel vastused lakkavad viitamast varem kokkulepitud tingimustele ja piirangutele \cite{datastudios2025perplexity}. Seetõttu ei sobi Perplexity lõputöö genereerimise tööriistaks, kus on vajalik suur ja järjepidev kontekst.

\textbf{GitHub Copilot Pro ei sobi akadeemilise teksti genereerimiseks.} GitHub Copilot Pro, mida Tartu Ülikool pakub üliõpilastele tasuta GitHub Education programmi kaudu, on mõeldud eelkõige koodi kirjutamise abivahendiks \cite{github2026copilot}. GitHubi enda dokumentatsiooni kohaselt on Copilot Chat mõeldud ainult kodeerimisega seotud päringutele ning juhiste piiramine kodeerimisülesannetega parandab mudeli väljundi kvaliteeti \cite{github2026copilot_responsible}. Lõputöö genereerimise kontekstis osutus see aga ebapiisavaks. Copilot'i mudelid tundusid oluliselt ``laisemad'' -- nad kippusid andma lühikesi, pealiskaudseid vastuseid ja vältisid pikemate akadeemiliste tekstilõikude genereerimist. Nguyen ja Nadi \cite{nguyen2022empirical} on oma empiirilises uurimuses näidanud, et Copilot'i soovitustel on tähemärkide pikkuse piirang, mis takistab pikema väljundi genereerimist, ning keerulisemate ülesannete korral langeb soovituste kvaliteet märgatavalt. Lisaks rakendab GitHub Copilot'ile päringupõhiseid limiite (\textit{rate limits}), mis piiravad kasutamise intensiivsust \cite{github2026copilot_ratelimits}. Kokkuvõttes on tööriist optimeeritud lühikeste koodilõikude ja kooditäienduste jaoks, mitte terviklike akadeemiliste tekstide loomiseks. Seetõttu ei saa GitHub Copilot Pro-d soovitada lõputöö kirjutamise abivahendina, vaatamata sellele, et see on üliõpilastele tasuta kättesaadav.


\subsection{Eetilised kaalutlused ja akadeemiline ausus} \label{subsec:eetika}

TI kasutamine lõputöö kirjutamisel tõstatab olulisi eetilisi küsimusi. Tartu Ülikooli akadeemilise aususe reeglid \cite{tuati2026juhend} nõuavad, et üliõpilane:
\begin{itemize}
    \item deklareerib kõik TI vahendid, mida ta kasutab;
    \item märgib, millises töö osas ja millises ulatuses TI-d kasutati;
    \item tagab, et töö peegeldab tema isiklikku akadeemilist arengut ja analüütilist mõtlemist.
\end{itemize}

Rodafinos \cite{rodafinos2025integration} rõhutab, et ülikoolid peaksid TI kasutamise keelustamise asemel lõimima TI-kirjaoskuse õppekavadesse ja kehtestama selged eetilise kasutamise poliitikad. Lund et al. \cite{lund2025perceptions} leidsid, et tudengid ei käsitle TI kasutamist kirjutamisel tavalise plagiaadi laiendusena, vaid eraldiseisva akadeemilise käitumise kategooriana, mis nõuab eraldi eetilist raamistikku.

On oluline rõhutada, et käesolev töö ei propageeri lõputöö TI abil ``kirjutamist'', vaid uurib TI kasutamist \textit{abivahendina} kirjutamisprotsessis. Nii nagu õigekirjakontroll ei kirjuta teksti üliõpilase eest, ei tohiks ka TI-põhine tekstigenereerimine asendada üliõpilase isiklikku panust. TI-d tuleks käsitleda kui tööriista, mis aitab struktureerida mõtteid, luua mustandeid ja parandada keelekasutust.

\subsection{Tulevikuväljavaated} \label{subsec:tulevikuvaljavaated}

TI mudelite areng on kiire ja käesoleva töö tulemused peegeldavad 2025--2026. aasta seisu. Lähitulevikus on oodata järgmisi arenguid:

\begin{itemize}
    \item \textbf{Eesti keele kvaliteedi paranemine} -- uuemate mudelite (nt oodatav GPT-6, Gemini 4.0) treeningandmetes on eesti keele osakaal tõenäoliselt suurem, mis parandab keelelist kvaliteeti. EstLLM projekti \cite{dorkin2026estllm} tulemused näitavad, et spetsialiseeritud jätkutreenimine võib eesti keele kvaliteeti märkimisväärselt tõsta ka olemasolevates mudelites.
    \item \textbf{Agentsed akadeemilised tööriistad} -- Claude Code'i, Cowork'i, Google'i Antigravity ja OpenAI ChatGPT Codex laienduse laadsed agentsed tööriistad, mis suudavad iseseisvalt faile hallata, kompileerida ja projekti struktureerida, muutuvad tõenäoliselt akadeemilise kirjutamise standardtööriistadeks. Meie eksperiment kinnitas, et agentne töövoog parandab oluliselt genereerimise kvaliteeti kõigi testitud mudelite puhul -- nii Opus 4.6 ja Sonnet 4.6 (Cowork), Gemini 3.1 Pro (Antigravity) kui ka GPT-5.4 (Codex) puhul.
    \item \textbf{Viidete kvaliteedi paranemine} -- RAG (Retrieval-Augmented Generation) süsteemide areng peaks parandama viidete täpsust ja vähendama hallutsinatsioone. Kresevic et al. \cite{kresevic2024rag} on näidanud, et RAG-raamistik koos struktureeritud konteksti ja promptide disainiga parandas GPT-4 Turbo täpsust 43\%-lt 99\%-le, mis demonstreerib selle lähenemise potentsiaali ka akadeemilises kontekstis. Ji et al. \cite{ji2024cot_citations} leidsid, et ahelpõhjendamise (\textit{chain-of-thought}) kasutamine parandas oluliselt viidete kvaliteeti teksti genereerimisel.
    \item \textbf{Kontekstiakende suurenemine} -- kontekstiaknad kasvavad jätkuvalt, mis võimaldab mudelitele sisestada üha rohkem infot korraga.
    \item \textbf{Spetsialiseeritud akadeemilised mudelid} -- lähiaastatel on tõenäoline spetsiaalselt akadeemiliseks kirjutamiseks kohandatud mudelite (nt eesti keelele ja Tartu Ülikooli stiilile peenhäälestatud variandid) ilmumine, mis võiksid ühendada üldmudelite võimekuse valdkonnaspetsiifiliste teadmistega ning vähendada hallutsinatsioonide osakaalu.
    \item \textbf{TI-genereeritud sisu tuvastamine ja integreerimine õppekavasse} -- ülikoolide ja kirjastuste poolt kasutusele võetavad TI-detekteerimise tööriistad arenevad paralleelselt genereerimistööriistadega, samas kui akadeemilised institutsioonid liiguvad TI keelustamiselt selle läbipaistva integreerimise suunas. See nõuab uusi hindamismetoodikaid, mis keskenduvad üliõpilase sisulisele panusele ja mõtteprotsessile, mitte ainult lõpptulemusele.
    \item \textbf{Multimodaalsed võimed akadeemilises tekstis} -- tulevased mudelid suudavad tõenäoliselt genereerida mitte ainult teksti, vaid ka akadeemiliselt kvaliteetseid jooniseid, diagramme ja visualiseerimisi otse \LaTeX{}-ühilduvas vormingus (nt TikZ või PGFPlots koodi), mis vähendab hetkel veel märkimisväärset käsitsitööd jooniste loomisel.
\end{itemize}

Käesolev uurimus piirdus neljast mudelist ja kolmest agentsest tööriistast koosneva võrdlusega 2025--2026. aasta seisuga. Arvestades TI valdkonna arengutempot, kus uusi mudeleid avaldatakse mitu korda aastas, tuleks sarnaseid teostatavusuuringuid korrata regulaarselt, et hinnata, kuidas ülalnimetatud arengud tegelikult realiseeruvad ja millised käesoleva töö järeldustest jäävad kehtima ka uuema põlvkonna mudelite puhul.

